\newcommand{\DoNotLoadEpstopdf}{}
\documentclass[11pt,reqno]{amsart}
\usepackage[T1]{fontenc}
\usepackage{lmodern,amsmath,amssymb,amsthm,mathtools,mathrsfs,microtype}
\usepackage[margin=1.1in]{geometry}
\usepackage{xr-hyper}
\usepackage[hidelinks]{hyperref}
\externaldocument[supp-]{build/supplement}[supplement.pdf]
\hypersetup{pdftitle={Collision records, arithmetic, and raw local inversion in a triangular Lorentz gas},pdfauthor={Qian Qi}}
\setlength{\emergencystretch}{2em}
\allowdisplaybreaks[2]
\numberwithin{equation}{section}
\newtheorem{theorem}{Theorem}[section]
\newtheorem{lemma}[theorem]{Lemma}
\newtheorem{proposition}[theorem]{Proposition}
\newtheorem{corollary}[theorem]{Corollary}
\theoremstyle{definition}
\newtheorem{definition}[theorem]{Definition}
\newtheorem{inputstatement}[theorem]{Input statement}
\theoremstyle{remark}\newtheorem{remark}[theorem]{Remark}
\newcommand{\R}{\mathbb R}\newcommand{\Z}{\mathbb Z}
\newcommand{\one}{\mathbf1}\newcommand{\dd}{\,\mathrm d}
\newcommand{\esssup}{\operatorname*{ess\,sup}}
\newcommand{\diver}{\operatorname{div}}
\newcommand{\TV}{\operatorname{TV}}
\title[Collision records and raw local inversion]{Collision records, arithmetic, and raw local inversion in a triangular Lorentz gas}
\author{Qian Qi}
\date{October 11, 2026. A2-DYN, revision 80}
\subjclass[2020]{37D50, 37A50, 60F05, 28A75}
\keywords{Sinai billiards, exact collision records, local limits, singular insertions, homogeneity strips, optical drift, exact-label coarea flux}
\begin{document}
\raggedbottom
\begin{abstract}
We study the pointwise local limit problem for the complete displacement,
collision-count, occupation and roof record of actual returns in a
triangular Lorentz gas. Homogeneous grazing strips give a superlinear
weak estimate for real-frequency physical powers. Combining it with
a fixed refinement of the spectral peaks and a physical cap on the
remainder yields a marked-incidence bound which can be integrated
against inverse incidence. This proves a complete fixed-width local
estimate for unbounded insertions, without assuming that they are
bounded strong-space multipliers. Applied to the canonical optical
flow, it controls inverse-twist loss and the initial mass of genuine
grazing exits at every internal incidence scale, on sources with
noncritical initial endpoint momentum. The same drift estimate pays
the compression cost in the nongrazing first-exit pullback. Common
coarea kernels and a flow-box exhaustion keep all exact labels and
true boundaries visible. The original scalar endpoint is reduced to
the critical initial source and the remaining nongrazing boundary
flux; their collision-uniform height estimates are not inferred from
fixed-width volume concentration.
\end{abstract}
\maketitle

\section{The unchanged physical problem}
Put $R_-=9/20$ and $R_+=47/100$, and let $R\in I=[R_-,R_+]$. Circular scatterers have centers in
$\Lambda=\Z(1,0)+\Z(1/2,\sqrt3/2)$. Write $T_R$ for the regular
collision map and $(\alpha,p)\in(\R/2\pi\Z)\times(-1,1)$ for its
coordinates; $p$ is the outgoing tangential velocity and
$\dd\nu=\dd\alpha\,\dd p/(4\pi)$. Physical flights satisfy
$\tau_-\le\tau_R\le\tau_+$, where $\tau_-=3/50$ and $\tau_+$
is the fixed finite-horizon bound.

The section is specified here to fix every event in the argument.
Put $h_0=1/200$. Let $Y_R$ be the union of the four open coordinate
squares of half-side $h_0$ centered at
\[
 (0,0),\quad(\pi/6,-1/2),\quad
 (-\theta_R,\sin\theta_R),\quad
 (\pi/3-\theta_R,\sin\theta_R).
\]
Here $\theta_R=2\arctan q_R$, and $q_R$ is the unique root in
$(49/100,53/100)$ of
\[
 8Rq(1+q^2)-4q(1-q^2)+\sqrt3(1-6q^2+q^4)=0.
\]
Set
\[
 U=\{|\alpha-\pi/6|<3/50,\ |p|<3/20\},\qquad
 Y_R^*=Y_R\cup U,\qquad c_*=\frac{91}{10000\pi}.
\]
The rectangles are disjoint and $\nu(Y_R^*)=c_*$. We always use
$\nu_R^*=c_*^{-1}\nu|_{Y_R^*}$, not a renormalization of a smaller
source.

At $m$ flights, let $L=L_{m,R}$ be their total length,
$K^{\rm c}_{m,R}$ their lattice displacement, and
$A_{m,R}=\sum_{j=0}^{m-1}\one_{Y_R^*}\circ T_R^j$. The event is
\begin{equation}\label{eq:v76-event}
 \Pi_{n,k,m,R}=Y_R^*\cap\{K^{\rm c}_{m,R}=k,\ A_{m,R}=n\}
                         \cap T_R^{-m}Y_R^*.
\end{equation}
The initial visit is counted; the terminal visit is tested separately.
On this event the actual $n$th return has collision count $m$ and
roof $L$. Denote its density by $p_{n,R}(k,m,t)$.

Write $c_j=\sqrt{1-p_j^2}$ and let $\Delta_j$ be the least clearance
of flight $j$ from any third disk, excluding its two incident disks.
Keep the original smooth function $0\le\Theta\le1$, zero on
$(-\infty,1]$ and one on $[2,\infty)$. With $q_0=47/53$,
$d_j=\min(j,m-j)$ and $d'_j=\min(j,m-1-j)$, set
\begin{align}
 H^\varepsilon&=
 \prod_{j=0}^{m}\Theta\!\left(\frac{c_j}{\varepsilon q_0^{d_j/4}}\right)
 \prod_{j=0}^{m-1}\Theta\!\left(\frac{\Delta_j}{\varepsilon q_0^{d'_j/4}}\right),
 \qquad W^\varepsilon=1-H^\varepsilon,\label{eq:v76-guard}\\
 b^\varepsilon(t)\dd t
   &=L_\#(W^\varepsilon\one_{\Pi_{n,k,m,R}}\nu_R^*).
 \label{eq:v76-source}
\end{align}
In particular first incidence and first clearance outside any selected
critical neighborhoods are still present. No section-distance factor
has been inserted into this guard.

\subsection{The arithmetic reference and the target}
For $y=(k_1,k_2,t,n)$ put
\[
 \mu_R=(0,0,\bar\tau_R,c_*),\quad
 Z=(y-m\mu_R)/\sqrt m,\quad
 \mathcal C_{m,R,n,k}(M)=\{t:|Z|\le M\}.
\]
We retain the finite spectral-peak reference, including peaks close to
but not on the unit circle. To specify its meaning, the collision
record observable is $(\kappa_1,\kappa_2,\tau,\one_{Y_R^*})$.
Let $\lambda_j(R,\vartheta)$ and $\mathsf P_j(R,\vartheta)$ be its
continued simple twisted eigenvalue and projection in each of the
fixed peak charts, and let $\chi_j$ be their frequency partition.
At the real maximum $a_j$ of $\log|\lambda_j|$ put
\[
 z_j=\lambda_j(R,a_j),\quad
 \mu_j=i^{-1}\nabla\log\lambda_j(R,a_j),\quad
 H_j=-\nabla^2\log\lambda_j(R,a_j),
\]
with $\operatorname{Re}H_j\ge aI_4$ for a common $a>0$. If
$\eta_R=\one_{Y_R^*}$ and $\ell$ denotes mass evaluation, put
$B_j=\chi_j(R,a_j)\ell(M_{\eta_R}\mathsf P_j(R,a_j)(\eta_R\nu))$.
Define
\begin{align}
 G_H(z)&=(2\pi)^{-4}\int_{\R^4}e^{-iv\cdot z-v^{\mathsf T}Hv/2}\dd v,
 \notag\\
 \mathcal L_{m,R}(y)&=\operatorname{Re}\left[
 c_*^{-1}\sum_j e^{-ia_j\cdot y}z_j^m B_j
      G_{H_j}\left((y-m\mu_j)/\sqrt m\right)\right].
 \label{eq:v76-reference}
\end{align}
Only the charts active at $R$ are summed. This is the same transition
reference as in the retained spectral theorem. The existence and
uniform estimates for these scalar data belong to that analytic
input, not to the transport argument below. No frequency branch or
arithmetic zero class is deleted.

The governing assertion remains
\begin{equation}\label{eq:v76-target}
 \sup_{R,\,1\le n\le m,\,k\in\Z^2}
 \esssup_{t\in\mathcal C_{m,R,n,k}(M)}
 |m^2p_{n,R}(k,m,t)-\mathcal L_{m,R}(k_1,k_2,t,n)|\longrightarrow0.
\end{equation}
Let $\mathcal E_M$ denote this error's collision limsup. An empty
central target contributes zero. The order of the uniformity is
$\sup_R\esssup_t$, so exceptional density sets may depend on $R$.

\subsection{The analytic inputs in the form used here}
\begin{inputstatement}\label{input:v76-analytic}
The unchanged source \eqref{eq:v76-source} has a nonnegative,
compactly supported $BV$ density at every fixed count. For every
fixed $\varepsilon,\delta>0$, with
$A_\delta f(t)=(2\delta)^{-1}\int_{t-\delta}^{t+\delta}f(s)\dd s$,
\begin{equation}\label{eq:v76-average}
 \limsup_m\sup_{R,n,k}m^2\|A_\delta b^\varepsilon\|_\infty
                                      \le C\varepsilon^{1/16}.
\end{equation}
The constant is independent of the previously fixed width. For large
fixed $B$ and $\varepsilon(B)=A_0B^{-1/12}$,
\begin{equation}\label{eq:v76-remainder}
 m^2p-\mathcal L=m^2b^{\varepsilon(B)}+e_B,\qquad
 \limsup_m\sup_{R,n,k}\esssup_{\mathcal C(M)}|e_B|
                                     \le C_MB^{-1/192}.
\end{equation}
All auxiliary spectral bands in these statements are fixed before
the collision limit and enlarged only afterwards.
\end{inputstatement}
These are the complete-source tame theorem, the fixed-width
complete-source bound, and the positive scalar remainder theorem in
the technical supplement. Their proofs and historical hypotheses are
retained there. The fixed-count source uses threshold integration
\cite{CMInt}, definable monotonicity \cite{Coste}, and the
distributional characterization of $BV$ \cite{AFP}; it does not
assume definability of the original smooth guard. They are restated as inputs to distinguish them from
the new geometric and transport proofs. No long-time bound on a
fixed-count variation constant is part of this statement.



\subsection{Precise inherited inputs}\label{sec:v79-input-map}
The fixed-width estimate is Supplement Theorem~\ref{supp-thm:v49-physical-remainder-local},
in particular Supplement equation~\eqref{supp-eq:v49-local-limsup}.
It uses Supplement Theorem~\ref{supp-thm:v49-physical-marked-local}, on the
original exact displacement and occupation coefficient. For an interval
of length $2r$, division by $2r$ gives \eqref{eq:v76-average}. That
proof fixes the auxiliary band, takes the collision limit, and only
then enlarges the auxiliary band. Its leading constant is independent
of the previously fixed interval length. The extension from small
margins to a fixed remaining compact range is stated there explicitly.
At auxiliary parameter one, $\Theta(c_0)=0$ since $c_0\le1$, so $W^1=1$
and $b^1=p$. This supplies the unweighted fixed-width estimate without
asserting a raw pointwise bound. The positive remainder is the inherited
scalar input \eqref{eq:v76-remainder}, stated as Supplement
Theorem~\ref{supp-thm:v51-positive-error}; no surface-current estimate is
included in any of these volume-source inputs.

\subsection{Singular insertions and physical optical transport}
The first-variation coordinates are $q=(\alpha_0,\alpha_m)$,
$g=\nabla_qF=R(-p_0,p_m)$, $K=D_q^2F$ and
$\vartheta=\arg(-p_0,p_m)$. The canonical field satisfies
$XF=1$, $X\vartheta=0$, and is stopped at the true boundary of the
regular exact event. Put
$\omega=|\partial_{p_0}\alpha_m|^{-1}$ and
$\chi=\det D_q(F,\vartheta)$. The retained geometric proofs below
give uniform positive bounds for $K$ and $\chi$, a uniform cubic
bound for $F$, and the physical Jacobian factorization
\[
             J_h=\frac{\omega(\Phi_hx)}{\omega(x)}
                       \frac{\chi(x)}{\chi(\Phi_hx)}.
\]
The angular ratio and the original guard have already been controlled
on the complete source. The remaining optical sensitivity contains
$\sum_i q_0^{d_i}/c_i$, which is singular at small internal incidences.

The new estimate for this sensitivity is a local concentration theorem
for the actual unbounded volume insertion
\[
 I_m=1+\sum_{i=1}^{m-1}q_0^{d_i}c_i^{-1},\qquad
 u(t)\,dt=L_\#(I_m\one_{\Pi_{n,k,m,R}}\nu_R^*).
\]
Corollary~\ref{cor:v80-depth-source} proves, for each fixed $r>0$,
\[
              \limsup_m\sup_{R,n,k}m^2\|A_r u\|_\infty\le C.
\]
The proof uses a capped superlinear marked-incidence estimate, not a
pointwise bound for another bad source and not a bounded multiplier
assertion for $c^{-1}$. All grazing thresholds are summed before the
collision limit, while the auxiliary Fourier band remains fixed.

Let $E_h^\varepsilon$ be the original weighted true-exit mass, and
let $O_h^\varepsilon$ be the positive inverse-twist loss on survivors,
as defined in Section~\ref{sec:v79-volume}. For a fixed
$a>0$, restrict their initial source to $|g|\ge a$; use $\mathrm{gr}$
for a true exit with some internal incidence tending to zero.
Theorem~\ref{thm:v80-noncritical-optical} gives the full estimate
\[
 \limsup_m\sup_{R,n,k}m^2
       \|\bar O_r^{\varepsilon,a}
           +\bar E_r^{\varepsilon,\mathrm{gr},a}\|_\infty
                  \le Cr/a.
\]
This includes every positive internal incidence, without an incidence
corridor. The zero optical density at a grazing exit does not erase
its positive starting mass: that mass is paid by the negative optical
variation. The corresponding nongrazing boundary inequality is
Theorem~\ref{thm:v80-boundary-reduction}; it retains the outgoing
physical flux and explicitly pays the missing pullback by volume drift.

Corollary~\ref{cor:v80-endpoint-budget} inserts these estimates into
the original pointwise objective. The initial critical neighborhood
and the nongrazing boundary flux still require their own height
estimates. An integrable singular insertion at every fixed width is
not a uniform pointwise bound for a critical Abel convolution, nor
is it a theorem about a surface-current insertion.

\subsection{Contact estimates and the literal guard}
The constant in each entry below is uniform in the collision count.
Here $\ell_q$ is the length of a regular endpoint curve and the displayed
powers refer to the contact and flight depths defined above.
\begin{center}
\begin{tabular}{lll}
Quantity & Change bounded by & Original guard factor\\[2pt]
\hline
Incidence $c_j$ & $Cq_0^{d_j}\ell_q$ & Yes\\
Third-disk clearance $\Delta_j$ & $Cq_0^{d'_j}\ell_q$ & Yes\\
Section-edge distance & $Cq_0^{d_j}\ell_q$ & No\\
Flight length $\ell_j$ & $Cq_0^{d'_j}\ell_q$ & No
\end{tabular}
\end{center}
Section decisions remain exact event tests. They are not new smooth
guard factors. The canonical field crosses coordinate cuts by change
of chart; it stops at the prescribed physical and section-decision
boundaries. Exhaustion boundaries enter only the limiting flux identity
in Proposition~\ref{prop:v79-source-exit-flux}.

\input{core/174_canonical_endpoint_momentum_flow}
\input{core/177_all_contact_motion_and_guard}
\input{core/180_uniform_cubic_reduced_action}
\input{core/181_separated_physical_volume_budget}
\input{core/182_optical_sensitivity_and_exit_flux}

\input{core/183_homogeneous_incidence_localization}
\input{core/184_singular_optical_transport}
\input{core/185_common_coarea_and_remaining_boundary}

\section{The remaining scalar estimates and their consequences}
The organizing assertion remains \eqref{eq:v76-target}, with the
complete arithmetic transition, original normalization, exact labels
and original smooth guard. The new singular insertion theorem closes
the incidence-scale summation needed for optical drift away from the
initial critical set. The boundary inequality estimates the same
pullback cost for nongrazing exits, but does not bound their outgoing
flux height. The two residual sources are displayed in
\eqref{eq:v80-endpoint-budget}; neither is suppressed by changing the
event or by a signed occupation-current cancellation.

All widths, including the initial critical threshold, are fixed
before the collision limit. After that limit the auxiliary spectral
band is enlarged, and the outer reconstruction parameter is taken
last. The full scalar theorem is not asserted until the two remaining
height estimates have been proved in this order. The new capped
remainder, common-kernel argument and singular-volume estimate do not
license exchanging these limits.

After scalar closure, Supplement
Theorem~\ref{supp-thm:v73-scalar-path-equality} and the positive path
remainder give the same-roof path numerator. Normalization is made
only on positive arithmetic reference classes. Collision-to-return
clock transfer and forward likelihood retain their separate inherited
hypotheses; arithmetic zero classes are never divided by. The complete
preceding article and every inherited proof remain in the supplement.

\input{journal-references}
\end{document}
